%! Comparing Distributions — Unit 2, Lesson 2.4 — Lesson Plan
% GRADUAL-RELEASE plan, Unit 2 timing 5 / 45 / 5 / 5 (user decision 2026-09-29): warm-up /
% guided notes and practice / spoken debrief / close and ASSIGN the homework (not started in
% class). The guided notes carry the release ROW BY ROW in one Main Ideas / Notes table: I do /
% we do in rows 1-4 (problems 1-16, ~25 min) -> we do the Guided Practice row (problems 17-20,
% ~20 min). THERE IS NO INDEPENDENT PRACTICE SET IN THE NOTES — the homework is the individual
% practice. Teacher-facing. Breakable boxes are `skillbox` with a \boxguard ahead; the two
% tabularx boxes are `fixedskillbox` (never splits). No group activity, no exit ticket.
\documentclass[10pt]{article}
\usepackage{saar-article}
\usepackage{saar-boxes}
\usepackage{graphicx}   % -article does NOT load graphicx
\graphicspath{{images/}}
\newcommand{\TallMath}[1]{$\displaystyle #1\rule[-1.4em]{0pt}{3.2em}$}
\newcommand{\UnitNumberName}{Unit 2: Describing One Variable \quad}
\newcommand{\LessonNumberName}{Lesson 2.4: Comparing Distributions}

\begin{document}
\begin{center}
    {\Large\bfseries \CourseName \\
    {\small\bfseries \UnitNumberName \LessonNumberName}}
\end{center}

\begin{tcolorbox}[colback=frost, colframe=cerulean, boxrule=0.9pt, arc=2mm,
    left=2mm, right=2mm, top=1.5mm, bottom=1.5mm]
    \textbf{Primary Objective:} Students will read two groups off a back-to-back stemplot,
    parallel dot plots, parallel boxplots, and calculator output on one common scale; write a
    comparison \emph{in context} that uses comparative words for shape, outliers, center, and
    spread; choose median and IQR or mean and SD to fit the shape; and justify which group is
    more consistent when the two centers are the same.
    \\[2pt]
    \textbf{Standards:} PS.DS.3a, PS.DS.3b \quad$\cdot$\quad carries forward Lessons 2.1--2.3
    (displays and SOCS; mean, median, SD, IQR; boxplots, $1.5\cdot$IQR fences, resistance)
    \quad$\cdot$\quad the last lesson before the Unit~2 review and test
    \\[2pt]
    \textbf{Lesson model:} \emph{Gradual release --- I do, we do, you do.} There is no group
    activity. The notes name the vocabulary as they teach it and deliver the instruction
    \textbf{row by row in one Main Ideas / Notes table} (four rows, problems 1--16: you read
    the definition, mark up the display, and work the first problem of each grid; the class
    works the rest); the Guided Practice row (problems 17--20) is worked together with students
    holding the pen. \textbf{The homework is the individual practice} --- assigned at Close \&
    Assign and done outside class, not started in class. The debrief is a \emph{spoken} cold
    check and one board item; there is no exit ticket.
    \\[2pt]
    \textbf{Conventions on the page:} quartiles are the medians of each half, leaving out the
    overall median when $n$ is odd (the Lesson~2.3 method, restated in row~2); Sx on the
    calculator screens is the \emph{sample} standard deviation (divide by $n-1$).
\end{tcolorbox}

\renewcommand{\arraystretch}{1.4}
\boxguard
\begin{skillbox}[Learning Targets \& Key Understandings]{goldbox}
\begin{tabularx}{\linewidth}{@{}X|X@{}}
\textbf{Learning targets (student language)} & \textbf{Key understandings (the ``why'')} \\[2pt]
\hline
\vspace{-6pt}
\begin{itemize}[leftmargin=1.1em, nosep, after=\vspace{2pt}]
  \item I can read a \textbf{back-to-back stemplot}, \textbf{parallel dot plots}, and
        \textbf{parallel boxplots}, and I compare groups only on a \textbf{common scale}.
        \hfill \textit{PS.DS.3a}
  \item I can read two groups' calculator output (1-Var Stats) and use it to compare them.
        \hfill \textit{PS.DS.3a}
  \item I can write a \textbf{comparison} in context with comparative words for shape,
        outliers, center, and spread. \hfill \textit{PS.DS.3b}
  \item I can choose median and IQR (skewed or outliers) or mean and SD (roughly symmetric,
        no outliers) to fit the shape. \hfill \textit{PS.DS.3b}
  \item I can decide which group is more \textbf{consistent} when the centers match, and
        explain why the spread decides. \hfill \textit{PS.DS.3b}
\end{itemize}
&
\vspace{-6pt}
\begin{itemize}[leftmargin=1.1em, nosep, after=\vspace{2pt}]
  \item \textbf{Target misconception:} that comparing distributions means comparing
        \emph{centers only} --- ``same median, so it doesn't matter'' --- or writing two
        separate descriptions side by side with no comparative word. \textbf{The rule that
        replaces it:} compare all four features with comparative words, in context; when the
        centers match, \emph{the spread decides}, and the group with less spread is more
        consistent.
  \item \textbf{Second misconception:} comparing two displays drawn on \emph{different}
        scales. A box that looks wider on a stretched axis can hold fewer minutes. Parallel
        displays exist so that both groups share one scale (PS.DS.3a).
  \item The measures follow the shape: an outlier or a long tail pulls the mean and SD, so
        median and IQR carry the comparison; mean and SD are the right pair only when both
        groups are roughly symmetric with no outliers.
  \item ``Better'' depends on the question --- \emph{within 40 minutes} favors the consistent
        shop, \emph{under 25 minutes} favors the spread-out one --- but either answer is read
        off the spread, never off the center.
\end{itemize}
\end{tabularx}
\end{skillbox}

\boxguard[24]
\begin{skillbox}[{Vocabulary, Concepts, \& Theorems --- taught directly in the guided notes}]{greenbox}
\small
\renewcommand{\arraystretch}{1.35}
\begin{tabularx}{\linewidth}{@{} >{\bfseries}p{3.4cm} X @{}}
Back-to-back stemplot & Two groups share one column of stems; one group's leaves go left, the
       other's go right. Every leaf is read with its stem, on either side. \\
Parallel dot plots / boxplots & Dot plots or boxplots of two or more groups stacked over one
       common scale, so they can be compared directly. \\
Common scale & One number line for every group, so equal distances mean equal amounts. \\
Comparison & A statement that uses comparative words (\emph{greater than, less than, about the
       same as}) for shape, outliers, center, and spread, in context. Two separate
       descriptions are not a comparison. \\
Consistent & Having less spread (smaller IQR or SD): the values stay close together. \\
Resistant measures & Median and IQR --- not pulled by outliers or a long tail; use them when
       either group is skewed or has an outlier (Lesson~2.3). \\
Sample standard deviation & Sx on the calculator --- the typical distance of values from the
       mean, dividing by $n-1$ (Lesson~2.2). \\
\end{tabularx}
\end{skillbox}

% fixedskillbox never splits, so the phase table stays intact on one page.
% THE PHASE TABLE IS A CONTRACT: 5 / 45 / 5 / 5 (Unit 2 on; 2026-09-29) — I do ~25 of the 45,
% Guided Practice ~20; the debrief is the spoken cold check; the homework is assigned, not started.
\begin{fixedskillbox}[Lesson at a Glance --- \MeetingLength]{frost}
\renewcommand{\arraystretch}{1.25}
\begin{tabularx}{\linewidth}{@{}l c X X@{}}
\textbf{Phase} & \textbf{Min} & \textbf{Students} & \textbf{Teacher} \\
\hline
Warm-Up & 5 & Tony's quartiles, IQR, and fence; two tiny lists with the same mean --- silent &
  Debrief item~3 aloud; write ``same mean, different spread'' on the board and leave it up all
  period \\
Guided Notes \& Practice & 45 & Fill the vocabulary box as each term is named; work rows 1--4
  (problems 1--16) with the pen in their hand; then the Guided Practice row, \emph{Paper
  Airplanes} (17--20), together & \textbf{I do / we do, row by row}: read the definition, mark
  up the display, work problems 1, 5, 10, 14; the class works the rest (25) $\to$ \textbf{we
  do} Guided Practice 17--20 (20) \\
Debrief & 5 & Pens down --- answer aloud, nothing to fill in & Put problem~19 back on the
  board (8 of 10 vs.\ 4 of 10); ask the cold check \\
Close \& Assign & 5 & Write down the homework, the due date, and item~6 as the first item &
  Announce the homework and its form (packet), preview the Unit~2 review \\
\end{tabularx}
\end{fixedskillbox}

\begin{fixedskillbox}[Warm-Up --- Activate Prior Knowledge (5 min)]{frost}
\begin{tabularx}{\linewidth}{@{}X X@{}}
\begin{minipage}[t]{\linewidth}\vspace{0pt}
    \textbf{The three items and what each seeds}
    \begin{itemize}[leftmargin=1.1em, itemsep=1pt, topsep=2pt]
      \item \textbf{1.} Tony's $Q_1 = 22$, median $25$, $Q_3 = 30$, IQR $8$. \textbf{Spirals}
            to Lesson~2.3 --- and they are the numbers row~1's stemplot and row~2's boxplot
            are built on (problems 2 and~6).
      \item \textbf{2.} Upper fence $30 + 1.5(8) = 42$, so $44$ is an outlier. \textbf{Seeds:}
            the dot at $44$ in row~2's boxplot, problem~7's lower-fence check on Rosa's, and
            homework item~2.
      \item \textbf{3.} Lists $9, 10, 11$ and $2, 10, 18$: both means $10$, ranges $2$ and
            $16$. \textbf{Seeds:} the crux --- row~4 (\emph{Different Spread}) is this
            sentence on real data.
    \end{itemize}
\end{minipage}
&
\begin{minipage}[t]{\linewidth}\vspace{0pt}
    \textbf{Running it}
    \begin{itemize}[leftmargin=1.1em, itemsep=1pt, topsep=2pt]
      \item \textbf{Debrief item~3 aloud} and write a student's sentence on the board ---
            \emph{same mean, different spread}. Leave it up all period; row~4 points back to
            it.
      \item Item~3's sentence is the tell. A student who writes \emph{``they are the
            same''} because the means match has walked in with today's misconception ---
            note who, and cold-call them at problem~14.
      \item Items 1--2 are Lesson~2.3's arithmetic --- say so, so nobody thinks they are
            being retested. Do not let them expand.
    \end{itemize}
\end{minipage}
\end{tabularx}
\end{fixedskillbox}

\boxguard
\begin{skillbox}[Hook --- before anyone sits down]{frost}
The hook is on \textbf{page 1 of the notes}, under the vocabulary box, and on the board:
\textit{``Rosa's Pizza and Slice House both advertise a median delivery time of $35$
minutes. Your party starts in $40$ minutes.''} \textbf{Ask: which shop do you call?}
Students circle \emph{Rosa's}, \emph{Slice House}, or \emph{it doesn't matter --- same
median} and write one line of reason \emph{before} anyone speaks; then take the hand vote and
write the split on the board. Most rooms vote \emph{it doesn't matter}. \textbf{Do not settle
it.} Say only that the answer is problem~14, in the \emph{Different Spread} row, and that they
will vote again there. The vote is what makes the crux land: students have to have committed
to \emph{same center, same shop} before the notes take it away.
\end{skillbox}

\boxguard
\begin{skillbox}[Guided Notes \& Practice --- I do, then we do (45 min)]{frost}
\textbf{This is the whole lesson.} There is no group activity, and there is no independent
practice set on this page: \textbf{the homework is the individual practice}, assigned at the
close and done outside class. The notes are \textbf{one Main Ideas / Notes table}: four
instruction rows (problems 1--16), then the Guided Practice row (17--20). \textbf{The rhythm
in every row is the same}: read the definition aloud, mark up the display, work the first
problem of the grid yourself, then hand the rest to the room. The vocabulary box on page~1 is
filled \emph{as each term is named}. The only blanks are the eight cells of the comparison
table (problem~9); every other answer is written in a problem's own space.
\begin{multicols}{2}
\textbf{I do / we do, row by row (about 25 min).}

\emph{Row 1, Back-to-Back Stemplots (problems 1--4).} Tony's warm-up times are the left
side. Name the \emph{back-to-back stemplot}, read the key together ($25$ and $24$), then work
problem~1 yourself (fastest and slowest). The class does 2--4. \textbf{Problem~3 is the
trap}: the left leaves read outward, and a student who says $81$ has read leaf-then-stem.
Five minutes.

\emph{Row 2, Parallel Displays (problems 5--8).} The same two shops as parallel dot plots and
parallel boxplots over one axis; students label the two pairs. Name \emph{common scale} as you
run a finger up a gridline. Work problem~5 yourself (Rosa's box sits to the right --- longer
deliveries). The class does 6--7. \textbf{Problem~8 is the second trap}: Mia's separate
scales make Rosa's box look wider when its IQR is smaller. Six minutes.

\emph{Row 3, Compare in Context (problems 9--12).} Read the definition of a
\emph{comparison} and the measures-fit-the-shape sentence. Take problem~9's table from the
class quickly (the word bank carries it), then read the frame aloud and work problem~10
(centers) yourself. The class does 11--12; the \emph{why median and IQR} justification is
said aloud from the frame, and problem~17 makes students carry it. Six minutes.

\emph{Row 4, Different Spread (problems 13--16) --- the crux.} Back to the hook's two shops
on one scale, medians marked. Work problem~13 live (both medians $35$; IQRs $6$ and $23$),
then \textbf{re-take the hook vote at problem~14 before anyone counts}. Then count: Rosa's
$11$ of $11$ within $40$ minutes, Slice House $7$ of $11$. Land the boxed sentence word by
word: \emph{when the centers match, the spread decides}. Problem~15 is the counterweight
(under $25$ minutes favors Slice House --- still a spread answer); problem~16 is Dev's two
lists, rewritten as a comparison. Eight minutes.

\columnbreak
\textbf{We do (about 20 min).} The Guided Practice row, \emph{Paper Airplanes} (problems
17--20) --- a fresh context, symmetric data, and pre-made calculator output, so the
\emph{mean and SD} branch finally gets used. \textbf{Students hold the pen; you hold the
questions.} Problem~17 decides the measures (both symmetric, no outliers $\to$ mean and SD).
Problem~18 compares with comparative words (same mean $30$ ft; SD $7.26$ vs.\ $2.58$ ft ---
almost three times). \textbf{Problem~19 is the crux transferred}: Kai says the same mean means
it doesn't matter. Take a hand count on Kai before anyone counts the zone; then Design A
lands $8$ of $10$ in the $27$--$33$ ft zone, Design B $4$ of $10$. Problem~20 is the full
comparison sentence the homework will ask for.

Circulate during Guided Practice and demand the reason, not the word:
\begin{itemize}[leftmargin=1.1em, nosep]
  \item \textbf{Problem 17:} ``What would make you switch to median and IQR?'' (A tail or an
        outlier --- neither is here.)
  \item \textbf{Problem 18:} ``You wrote two numbers. Where is the comparing word?''
  \item \textbf{Problem 19:} ``The means match. So what does decide it?'' A student who sides
        with Kai has not left the hook vote.
  \item \textbf{Problem 20:} ``Read me your sentence. Could a coach act on it?''
\end{itemize}
While circulating, pick the one paper to put up in the debrief --- a problem~19 that counted
the zone and said \emph{spread decides}. \textbf{This circulation is your formative read} ---
the homework is not started in class.
\end{multicols}
\end{skillbox}

\boxguard[24]
\begin{skillbox}[Debrief --- whole class, spoken (5 min)]{frost}
Pens down, papers up. \textbf{This debrief is spoken} --- there is no box at the end of the
notes to fill in, so it lives entirely on the board and in the room.
\begin{multicols}{2}
\begin{enumerate}[leftmargin=1.3em, nosep, itemsep=3pt]
  \item \textbf{The one board item: problem~19.} Put the paper you picked up (Design A $8$ of
        $10$ in the zone, Design B $4$ of $10$) and make a student say the sentence: \emph{same
        center, and the spread decides --- the design with less spread is more consistent.}
        Point back at the warm-up's line on the board: \emph{same mean, different spread}.
\end{enumerate}
\columnbreak
\textbf{Check for understanding --- ask it cold, whole class.} \emph{``Two classes took the
same $20$-point quiz. Both class means are $15$. Class~1's standard deviation is $1.2$
points; Class~2's is $4.5$. Maya says the classes did the same. Is she right? Say it as one
comparison sentence.''} Hands down, thirty seconds of think time, then cold-call three
students. Correct answer: \textbf{no} --- the centers are the same ($15$), but Class~2's SD
($4.5$) is much greater than Class~1's ($1.2$), so Class~1's scores are far more consistent.
\emph{This replaces the exit ticket.}

\textbf{Formative read.} Sort the answers into three piles: \textbf{(a)} said \emph{no},
named the same center, and compared the spreads with a comparative word; \textbf{(b)} said
\emph{no} but gave two separate descriptions (``Class~1 is $1.2$, Class~2 is $4.5$'');
\textbf{(c)} agreed with Maya --- centers only. \textbf{If pile (c) runs past a third of the
room, open the Unit~2 review by re-running row~4 (problems 13--16)}, and say so plainly.

\textbf{Five minutes: item~1 and the cold check only.} Anything else waits for the review.
\end{multicols}
\end{skillbox}

\boxguard
\begin{skillbox}[Homework --- scored, assigned today, due after two study halls]{goldbox}
Six items on \textbf{Volt vs.\ Nova phone batteries} ($11$ phones each; both medians $12$
hours): \textbf{1} each brand's five-number summary and IQR, and Nova's IQR in context (the
core procedure); \textbf{2} is Volt's $6$-hour phone an outlier --- the lower fence $11 -
1.5(2) = 8$; \textbf{3} Volt's mean with and without the $6$ ($11.7$ vs.\ $12.3$ hours, the
median $12$ both ways) and which pair of measures fits --- \emph{the spiral item}, reaching
back to Lessons 2.2--2.3, and the justification; \textbf{4} and \textbf{5} are the
\textbf{contrast pair}: the same question with opposite needs --- at least $10$ hours favors
Volt ($10$ of $11$ vs.\ $8$ of $11$), at least $16$ hours favors Nova ($3$ of $11$ vs.\ $0$)
--- two brands with the \emph{same} median sending two buyers opposite ways; \textbf{6}
Jordan's ``same median, so it doesn't matter'' --- \emph{the crux head-on}, answered with a
written comparison.

\textbf{Assigned at Close \& Assign and done outside class.} \textbf{Scoring:} 12 points --- 2
per item. \textbf{Item~6 is where the misconception shows}: a student who agrees with Jordan,
or who writes ``Volt: IQR 2; Nova: IQR 8'' with no comparing word, has not yet made the move.
\textbf{Items 4--5 carry the second check} --- a student who answers both by the median has
compared centers only.

\textbf{Packet or Desmos --- decided here.} The packet pages are authored and are the default.
The items this lesson exists for are \emph{written} comparisons and a buyer's decision, which
the packet carries and a graphing activity does not, so \textbf{2.4 is a packet night}. Say so
aloud at Close \& Assign.

\textbf{Preview of the next class.} The Unit~2 review: every tool from 2.1--2.4 on one data
set, then the Unit~2 test. Item~6 is test-style --- say so.

\textbf{Connections \& Big Ideas.} \textit{Unit 2 core idea:} a distribution is described by
its shape, outliers, center, and spread, and two distributions are compared on all four.
\textit{Standards covered:} PS.DS.3a (back-to-back stemplots, parallel dot plots, parallel
boxplots, with technology), PS.DS.3b (compare shape, center, spread, and unusual features in
context). \textit{Spirals forward to:} the Unit~2 review and test, and every later lesson
that compares two groups --- treatment against control in an experiment, one year against
the next. \textit{Reaches back to:} Lesson~2.1's displays and SOCS, Lesson~2.2's mean,
median, SD, and IQR, and Lesson~2.3's boxplots, fences, and resistance.
\end{skillbox}

\boxguard
\begin{skillbox}[Watch For (while circulating)]{redbox}
\begin{itemize}[leftmargin=1.2em, nosep]
  \item \textbf{``Same center, so it doesn't matter''} (the hook, problems 14 and~19, HW~4--6
        --- the target misconception). Probe: \emph{``The medians match. So which of these
        shops gets your pizza there on time --- and what number told you?''}
  \item \textbf{Two descriptions, no comparison} (problems 10--12, 16, 18, 20; HW~6):
        ``Tony's median is 25. Rosa's median is 35.'' Probe: \emph{``Put one word between those
        two numbers that tells me which is bigger --- and what that means for a customer.''}
  \item \textbf{Different scales} (problem~8): Mia's wider-looking box. Probe: \emph{``How many
        minutes does that box cover? Put both on one number line.''}
  \item \textbf{Leaf before stem on the left side} (problem~3): $81$ for $18$. Probe:
        \emph{``Which digit is the stem? Read it first, on both sides.''}
  \item \textbf{Mean and SD with an outlier} (row~3's frame, problem~17, HW~3): reporting Volt's mean of
        $11.7$ as its center. Probe: \emph{``What happened to the mean when the $6$ left?
        What happened to the median?''}
  \item \textbf{A comparison with no context} (problems 18 and~20): ``B's SD is bigger.''
        Probe: \emph{``Bigger SD of what, in what units --- and what does it mean for the
        contest?''}
\end{itemize}
Cold-call prompts: \emph{``Give me two groups at this school with about the same center and
very different spreads.''} \quad \emph{``When is the group with MORE spread the better
choice?''}
\end{skillbox}

\boxguard
\begin{skillbox}[Close \& Assign (5 min)]{goldbox}
\textbf{The homework is assigned here and done outside class --- it is not started in
class.} Hand the launch: six items on Volt vs.\ Nova phone batteries, packet pages, scored out
of 12 (2 per item), due at the start of the first class after two study halls. Point out that
the \emph{Keep in Mind} box on the cover and rows 3--4 of the notes are the pages to flip back
to. \textbf{Tell students to do item~6 first if they are short on time} --- it is the item
that shows whether row~4 landed --- and that items 4 and~5 are a pair: read both before
answering either. Then close on what changed: \emph{``You walked in knowing how to find a
median, an IQR, and a fence for one group. You walk out comparing two --- on one scale, with
comparative words, on all four features --- and knowing that when two centers match, the
spread decides.''} \textbf{Preview:} the Unit~2 review --- every tool from this unit on one
data set --- then the test.
\end{skillbox}

% ── Teacher notes — one per component, in packet order (three) ──────────────
% Teacher-only prose lives HERE, never in a _key. No note for the debrief — it is fully
% specified in its own box above.
\begin{teachernote}[Warm-Up]
Five minutes, silent, timed. Items 1 and~2 are Lesson~2.3's arithmetic on the same $11$ Tony's
times the notes open with --- say so, so nobody thinks they are being retested, and so row~1
can move fast (students already know Tony's median is $25$). \textbf{Item~3 is the one to read
while collecting}: students who write \emph{``same mean, but B is more spread out''} have made
today's distinction without the vocabulary. Students who write \emph{``they are the same''}
because the means match have walked in with the target misconception --- note two of them and
cold-call them at problem~14. Write a student's sentence on the board --- \emph{same mean,
different spread} --- and \textbf{leave it up all period}; the debrief points back at it.
\end{teachernote}

\begin{teachernote}[Guided Notes \& Practice]
\textbf{This one document is the lesson --- pace it 25 / 20, then 5 for the debrief and 5 to
assign the homework.} It is one table, and every row runs the same way: you read the
definition, mark up the display, and work the first problem of the grid (1, 5, 10, 14); the
class works the rest with the pen in their hand. Rows 1 and~2 must move --- five and six
minutes. Tony's and Rosa's numbers are the same across rows 1--3, so students are reading one
data set three ways, not learning three data sets.

\textbf{Spend the time in row~4, \emph{Different Spread}}, which carries the misconception.
Work problem~13 live, then \textbf{re-take the hook vote at problem~14 before anyone counts.}
If students count first, the vote is empty; if they commit to \emph{it doesn't matter} and
\emph{then} count $11$ of $11$ against $7$ of $11$, the rule arrives as the answer to a
question they already have. Problem~15 is the counterweight and keeps the rule honest: more
spread is sometimes what you want --- and it is \emph{still} the spread that answers.
\textbf{Do not skip problem~16}: Dev's two lists are the exact form the misconception takes in
writing, and the rewrite is the sentence the homework and the test ask for.

\textbf{Guided Practice (17--20) is where the pen changes hands.} \emph{Paper Airplanes} uses
symmetric data and calculator output, so it is the first place the mean-and-SD branch of row~3
is used. Problem~19 is the crux transferred: take a hand count on Kai before anyone counts the
zone. If time is short, protect 17, 19, and 20 --- problem~18's comparison can be said aloud.

\textbf{There is no independent practice set on this page, and that is deliberate --- the
homework is the individual practice}, assigned at the close and done outside class. Your
formative read is the Guided Practice circulation plus the cold check; if the centers-only
pile is more than a third of the room, \textbf{open the Unit~2 review by re-running
problems 13--16}, and say so plainly.
\end{teachernote}

\begin{teachernote}[Homework]
\textbf{Scored, 12 points (2 per item), due at the start of the first class after two study
halls; assigned at the close, not started in class.} Score items 1--3 for correctness (item~1:
all twelve cells right for 2, most right for 1; item~2 needs the fence \emph{and} the verdict;
item~3 needs both means and a reason naming the outlier). Score items 4--6 for the
\emph{reason}: in 4 and~5 the counts plus the brand; in 6, full credit only for a sentence that
names the same center \emph{and} compares the spreads with a comparative word in context.
\textbf{Item~6 is the one to read first.} Sort it into three piles --- compared the spreads in
context (ready); named the spreads but as two separate descriptions (ready, needs the
language); agreed with Jordan or compared medians only (the misconception survived). If more
than a third land in the last pile, open the Unit~2 review with problems 13--16 before
anything else. Items 4 and~5 are a diagnostic pair: a student who picks the same brand for
both buyers is choosing by center.

\textbf{Packet is the call for 2.4.} The review opens from item~6 --- put two anonymous
answers from different piles side by side on the board.
\end{teachernote}
\end{document}
